\documentclass[11pt]{article}
\usepackage[margin=1in]{geometry}
\usepackage{amsmath,amssymb,amsthm}
\usepackage[hidelinks]{hyperref}
\newtheorem{theorem}{Theorem}
\newtheorem{lemma}{Lemma}
\newcommand{\tr}{\operatorname{tr}}
\begin{document}
\begin{center}
{\Large A proof of the Factor 4 conjecture for Markov polynomials}\\[0.5em]
MathDB 366806\\
September 12, 2026
\end{center}

\paragraph{Task 1.}
Prove Conjecture 8.1 of Evans, Veselov, and Winn \cite{EVW}:
the coefficient at every integer point strictly inside the critical
triangle of a Markov polynomial is divisible by $4$.
Write its homogeneous numerator as
\[
 P_{a/b}(u,v,w)=
 \sum_{i,j} A_{ij}u^iv^jw^{a+b-1-i-j}.
\]
The critical triangle is given by $i<a$, $j<b$, and $bi+aj>ab$.
We prove the stronger congruence
\[
 P_{a/b}(u,v,1)\in (4,u^a,v^b)\,\mathbb Z[[u,v]].
 \tag{1}
\]
Homogeneity ensures that setting $w=1$ preserves every coefficient
$A_{ij}$ separately.

\begin{theorem}
For every coprime $1\leq a\leq b$, every coefficient $A_{ij}$ with
$i<a$ and $j<b$ is divisible by $4$. In particular, the Factor 4
conjecture holds.
\end{theorem}

We use the following elementary matrix observation.

\begin{lemma}\label{overlap}
Let $R=(\mathbb Z/4\mathbb Z)[[u,v]]$.
If matrices $E,F$ over $R$ are idempotent and
$F\equiv I-E\pmod2$, then $\tr(EF)=0$ in $R$.
\end{lemma}
\begin{proof}
Write $F=I-E+2H$. Expanding $F^2=F$ gives
$2(H-EH-HE)=0$.
Reducing the expression in parentheses modulo $2$ and multiplying
on both sides by $\overline E$ gives
$\overline E\,\overline H\,\overline E=0$.
Thus every entry of $EHE$ is divisible by $2$.
By cyclicity of trace and $E^2=E$,
$\tr(EH)=\tr(EHE)$ is divisible by $2$. Consequently
$\tr(EF)=2\tr(EH)=0$ in $R$.
\end{proof}

\begin{proof}[Proof of the theorem]
Put $S=1+u+v$ and define
\[
 M=\begin{pmatrix}v&1\\uv&u+1\end{pmatrix},
 \qquad
 N=\begin{pmatrix}v+1&1\\uv&u\end{pmatrix}.
 \tag{2}
\]
The generalized Cohn-matrix formula
\cite[Theorem 10.1]{EVW} gives a word $C_{a/b}$ with $b$ letters
$M$ and $a$ letters $N$ such that
\[
 \tr(C_{a/b})=S\,P_{a/b}(u,v,1).
 \tag{3}
\]
To match conventions, specialize the source matrices to $z=1$,
put $u=x^2,v=y^2$, and conjugate by $\operatorname{diag}(x,1)$.
Its initial matrices become $M$ and $MN$.
The Farey mediant construction concatenates words, so their letter
counts are $(b,a)$.
The trace formula can also be checked directly from
\[
 \tr(X^2Y)=\tr(X)\tr(XY)-\det(X)\tr(Y):
\]
the determinants of these words are $u^av^b$, and the identity gives
the same numerator Vieta recursion as in the source, with the
same initial values.

We next show that the order of the letters in (3) is immaterial
modulo $4$ under the trace. Direct calculation gives
\[
 MN-NM=2K,\qquad
 K=\begin{pmatrix}0&-1\\uv&0\end{pmatrix},
 \qquad \tr K=\tr(MK)=0.
 \tag{4}
\]
Modulo $2$, one has $N=M+I$, so any word reduces to a polynomial
in $M$. The identity $M^2-SM+vI=0$ reduces such a polynomial to
the form $\alpha I+\beta M$.
For an adjacent interchange in an arbitrary word, (4) gives
\[
 \tr(XMNY)-\tr(XNMY)=2\tr(YXK)=0\pmod4,
\]
because $\tr(YXK)=0\pmod2$.
Successive adjacent interchanges therefore prove
\[
 \tr(C_{a/b})=\tr(M^bN^a)\quad\hbox{in }R.
 \tag{5}
\]

We now work entirely in $R$. There are unique formal series
$\mu,\nu\in(u,v)R$ satisfying
\[
 \mu^2-S\mu+v=0,\qquad
 \nu^2-S\nu+u=0.
 \tag{6}
\]
For example, iteration of $\mu\mapsto v/(S-\mu)$ starting at zero
stabilizes each coefficient, since it increases the order of a
difference; the same construction applies to $\nu$.
It also shows $\mu\in vR$ and $\nu\in uR$.
Set $\lambda=S-\mu$ and $\eta=S-\nu$.
The series $\lambda,\eta,\lambda-\mu,\eta-\nu$ are units.
The characteristic polynomials of $M,N$ are precisely the two
quadratics in (6). Thus
\[
 E=\frac{M-\mu I}{\lambda-\mu},\qquad
 F=\frac{N-\nu I}{\eta-\nu}
 \tag{7}
\]
are idempotents, with
\[
 M=\lambda E+\mu(I-E),\qquad
 N=\eta F+\nu(I-F).
 \tag{8}
\]

The two idempotents are complementary modulo $2$.
Indeed, $N=M+I$ modulo $2$, and substitution in (6) shows that
the small root of the second quadratic is
$\nu=\lambda+1$ modulo $2$; its unit root is $\eta=\mu+1$.
Uniqueness of the small root justifies these identifications.
Equation (7) then gives $F=I-E$ modulo $2$.
Lemma \ref{overlap} yields
\[
 \tr(EF)=0\quad\hbox{in }R.
 \tag{9}
\]

Let $J=(u^a,v^b)R$. Since $\mu\in vR$ and $\nu\in uR$,
equation (8) and idempotence imply, modulo $J$,
\[
 M^b=\lambda^bE,\qquad N^a=\eta^aF.
\]
Equations (5) and (9) now give
$\tr(C_{a/b})=0$ modulo $J$.
Finally $S$ is a unit in $R$, so (3) implies
$P_{a/b}(u,v,1)\in J$.
Membership in this monomial ideal says exactly that every coefficient
with $i<a$ and $j<b$ vanishes modulo $4$.
This proves (1) and the theorem.
\end{proof}

\paragraph{Task 2.}
Characterize the critical-triangle coefficients whose exact
$2$-adic valuation is $2$, and determine geometric criteria for
divisibility by higher powers of $2$.

\paragraph{Task 3.}
Compute the corrections to trace invariance modulo $8$ and $16$
and use them to obtain finer coefficient congruences.

\paragraph{Task 4.}
Construct a free action of a group of order $4$ on the perfect
matchings contributing to each critical-triangle coefficient.

\paragraph{Task 5.}
Combine the matrix description with valuations adapted to the
Markov sail to prove its conjectured coefficient duality.

\begin{thebibliography}{9}
\raggedright
\bibitem{EVW}
S. J. Evans, A. P. Veselov, and B. Winn,
\emph{Arithmetic and geometry of Markov polynomials},
arXiv:2501.14882v3 (2025); current 28-page journal-hosted manuscript,
Conjecture 8.1 and Theorem 10.1.
\url{https://armj.math.stonybrook.edu/Articles/250708-Winn/article.pdf}.
\bibitem{E}
S. J. Evans, \emph{Quantum numbers and Markov polynomials},
doctoral thesis, Loughborough University, December 2025;
deposited February 9, 2026. Conjecture 2.6.3 and Theorem 2.7.5.
\url{https://doi.org/10.26174/thesis.lboro.31250974.v1}.
\end{thebibliography}
\end{document}
